\documentclass[10pt,a4paper]{amsart}
\usepackage[margin=19mm]{geometry}
\usepackage{amsmath,amssymb,mathtools}
\usepackage[scr=rsfs,cal=euler]{mathalfa}
\usepackage{xfrac}
\usepackage[unicode,hidelinks,bookmarks=false]{hyperref}
\newcommand{\ie}{i.e.\ }
\newcommand{\cf}{cf.\ }
\newcommand{\resp}{respectively\ }
\newcommand{\Subsection}{\subsection}
\providecommand{\nobreakdash}{\nobreak\hbox{-}}
\setlength{\emergencystretch}{2em}
\multlinegap0pt
\newcommand{\BA}{{\mathbb{A}}}
\newcommand{\BC}{{\mathbb{C}}}
\newcommand{\BN}{{\mathbb{N}}}
\newcommand{\BZ}{{\mathbb{Z}}}
\newcommand{\CH}{{\mathcal H}}
\newcommand{\Fy}{{\mathfrak{y}}}
\newcommand{\Fz}{{\mathfrak{z}}}
\DeclareMathOperator{\Hilb}{Hilb}
\newcommand\Hom{\operatorname{Hom}}
\newtheorem{thm}{Theorem}
\newtheorem{lemma}{Lemma}
\newtheorem{corollary}{Corollary}
\newtheorem{definition}{Definition}
\newtheorem{remark}{Remark}
\title[Theorem A of arXiv:2508.21717]{A proof of the Brian\c{c}on-Iarrobino Conjecture in three dimensions: the upper bound for the tangent space (Theorem A)}
\author{Owen Mackenzie and Fatemeh Rezaee}
\date{Source: arXiv:2508.21717v1 (CC BY 4.0)}
\begin{document}
\maketitle
\noindent\emph{Source and licence.} A proof of the Brian\c{c}on-Iarrobino Conjecture in three dimensions. Version arXiv:2508.21717v1 (CC BY 4.0), distributed under the Creative Commons Attribution 4.0 International licence, http://creativecommons.org/licenses/by/4.0/, as recorded in the arXiv OAI licence metadata for this exact version and shown on the abstract page https://arxiv.org/abs/2508.21717v1. Full rights notice: licenses/arxiv-2508.21717-CC-BY-4.0.txt.\par\smallskip
\noindent\textit{Editorial scope.} Counted unit: Theorem A (Upper bound) of the paper (native source0.tex lines 700--711), the bound $T(I)\leq(2m_1+1)l-2\binom{m_1+2}{4}$ for a zero-dimensional Borel-fixed ideal of colength $l=\binom{k+2}{3}+\Delta$, delivered together with the authors' own complete proof of it (source0.tex lines 1085--1119) as the single primary proof body. No mathematics is written, repaired or re-derived: statement, proof and all retained context are the authors' own text line for line, in the authors' own notation; printed slips of the source are kept literal. Required context retained uncounted, in the paper's own order: (a) lines 666--695, the paper's Section 0.2 basic definitions -- the point $[I]$ of $\CH=\Hilb^l(\BA^3)$ attached to a zero-dimensional ideal, the tangent-space dimension $T(I)$, pure exponents, Definition 0.2 (maximal singularity), Definition 0.3 (Borel-fixedness) and Remark 0.4 -- which fix the meaning of every hypothesis of the counted statement; (b) lines 881--1032, the paper's Section 1.6 with the graded space $H^{J,J'}=\mathrm{Hom}(J,S/J')$, the sets $A^{J,J'}$, $B^{J,J'}$, the dimension $\mathrm{hom}$, the height and the function $t(J,J')$, Lemma 1.5 with its complete proof (including the two claims proved inside it), Corollary 1.9, Corollary 1.10, the definitions of height and of $t(I_0,I_1)$, Lemma 1.14 with its complete proof and Corollary 1.15; the counted proof uses Corollary 1.15 and the double-binomial identity of the same section, and those statements are stated there in terms of the objects just listed (the authors' phrase ``as above''), which is why the whole subsection is retained rather than the corollary alone; (c) lines 1071--1084, the paper's Section 2.1 preparation -- the sentence introducing the proof, Lemma 2.1 with its complete proof, Corollary 2.2, and the sentence naming the Ramkumar--Sammartano argument the proof follows. Omitted from this package and not redistributed: the paper's title page, abstract, Section 0.1 history, the Examples of Sections 1.6 and 2 (lines 849--876 and 1033--1065), the rest of the source, the authors' macros file and modified bibliography style (only the nine notation macros the retained spans use are re-stated verbatim from the authors' own macros: $\BA$, $\BC$, $\BN$, $\BZ$, $\CH$, $\Fy$, $\Fz$, $\Hilb$, $\Hom$), the seven figure files, and the paper's complete bibliography (only the single entry cited inside the retained spans, printed as [37], is transcribed). The wrapper numbers its own theorem environments, so the retained Lemma 2.1 prints as Lemma 3, Corollary 1.15 as Corollary 3 and Corollary 2.2 as Corollary 4 (the Section 1.6 corollaries print as Corollary 1, Corollary 2 and Corollary 3), while the counted statement prints as Theorem 1 although the paper prints Theorem A; the paper's printed identifiers are named in this scope and in the item metadata, and the retained proof's own optional heading necessarily follows the wrapper's numbering. No retained span contains a graphics-inclusion command: the retained spans contain no \textbackslash includegraphics, \textbackslash includepdf, \textbackslash input or \textbackslash include, so no image asset is shipped or declared, even though the paper's own source contains seven \textbackslash includegraphics calls in sections that are not retained.
\section*{Retained context: basic definitions (paper \S0.2)}
To a $0$-dimensional ideal\footnote{An ideal $I$ such that the quotient $\BC[x,y,z]/I$ is Artinian (zero-dimensional).} $I$ in $R=\BC[x,y,z]$, we associate the corresponding point $[I]$ in the Hilbert scheme $\CH=\Hilb^l(\BA^3)$. We denote the dimension of the tangent space to the Hilbert scheme $\CH$ at $[I]$  by $T(I)$.



For a $0$-dimensional Borel-fixed ideal $I=(x^{m_1},y^{m_2},z^{m_3},\text{mixed generators})$, the exponents $m_i$ are called \emph{pure exponents}.

\begin{definition}[Maximal singularity]
    By a point $[I]$ in $\CH=\Hilb^l(\BA^3)$ \emph{having the maximal singularity} we mean that $T(I)\geq T(J)$ for any $[J]\in\CH$. 
    
    By convention, we use \emph{maximum singularity} if there is only one option.


\end{definition}

The following definition is crucial.

\begin{definition}[Borel-fixedness]\label{Def:Borel-fixedness} An ideal $I$ is called a Borel-fixed ideal,
if for any generator $g$ of $I$, we have 
\begin{itemize}
    \item if $z|g$, then the quotients $\frac{xg}{z}, \frac{yg}{z} \in I$, 
    \item if $y|g$, then the quotient  $\frac{xg}{y} \in I$. 
    \end{itemize}
\end{definition}




\begin{remark}\label{rem:Borel-fixed only}
    Note that for any $l\in \BZ^{\geq 1}$, there exists a Borel-fixed ideal $I$ (Definition \ref{Def:Borel-fixedness}) of colength $l$ which corresponds to a point of $\Hilb^l(\BA^N)$  with the maximal singularity. So, in this article, we only focus on Borel-fixed ideals.
    \end{remark}

\section*{Counted claim: Theorem A (Upper bound), paper \S0.3}
\begin{thm}[Upper bound]\label{thm:upperbound}
    Let 
    \begin{align*}
        I=(x^{m_1},y^{m_2},z^{m_3},\text{mixed generators})
    \end{align*}
    be a 0-dimensional Borel-fixed ideal in $\BC[x,y,z]$ of colength $l=\binom{k+2}{3}+\Delta$, where $k\in\BZ^{\geq1}$ and $0\leq \Delta \leq \binom{k+2}{2}-{1}$. Then 
    \begin{align*}
        T(I)
        \leq 
    (2m_1+1)l-2\binom{m_1+2}{4}. 
    \end{align*}
\end{thm}

\section*{Retained context: the graded spaces and inequalities used by the proof (paper \S1.6)}
\begin{lemma}\label{Lem:inequalityL1L2}
    Let $J$ and $J'$ be two $0$-dimensional monomial ideals in $S=\BC[y,z]$ of colengths $l$ and $l'$, respectively. Then, we have
    \begin{align}\label{eq: Ineq}
       \mathrm{hom}(J,S/J')\leq l+l'.
    \end{align}
\end{lemma}


\begin{remark}Note that we can still apply the method of calculating $\mathrm{hom}(J,S/J)$ used in \cite{Ramkumar-Sammartano} to calculate $\mathrm{hom}(J,S/J')$. Thus, the dimension of $\mathrm{hom}(J,S/J')$ in Lemma \ref{Lem:inequalityL1L2} will be calculated by summing the number of bounded connected components of $(\tilde{J}+\alpha)\setminus \tilde{J'}$ across all $\alpha=(\alpha_y,\alpha_z)\in\BZ^2$.
\end{remark}
Before proving Lemma \ref{Lem:inequalityL1L2}, we have the following definition.


\begin{definition}\label{Def:AJ}
The space $H^{J,J'}:=\mathrm{Hom}(J,S/J')$ is graded by $\mathbb{Z}^2$ according to the value of $\alpha$ used. We divide $H^{J,J'}$ into two subspaces, $H^{J,J'}=H^{J,J'}_p\oplus H^{J,J'}_n$ as
\[H^{J,J'}_p:=\bigoplus_{\alpha\in\mathbb{Z}\times\mathbb{N}}H_\alpha,\]
\[H^{J,J'}_n:=\bigoplus_{\alpha\in\mathbb{Z}\times\mathbb{Z}^-}H_\alpha.\]
We will see that that $\dim(H^{J,J'}_p)=l'$ and $\dim(H^{J,J'}_n)\leq l$. 

Let
\begin{align*}   
A^{J,J'}:=\{(U,\alpha)\in\mathcal{P}(\mathbb{N}^2)\times\mathbb{Z}^2:U\text{ is a connected component of }(\tilde{J}+\alpha)\setminus \tilde{J'}\}.\end{align*}

Lowering the grading $n,p$ to this set,  we get $A^{J,J'}_p$ and $A^{J,J'}_n$ corresponding to bases of $H^{J,J'}_p$ and $H^{J,J'}_n$ as follows. 
\[A^{J,J'}_p:=\{(U,\alpha)\in\mathcal{P}(\mathbb{N}^2)\times\mathbb{Z}^2:U\text{ is a connected component of }(\tilde{J}+\alpha)\setminus \tilde{J'},\alpha_z\geq 0\},\]
\[A^{J,J'}_n:=\{(U,\alpha)\in\mathcal{P}(\mathbb{N}^2)\times\mathbb{Z}^2:U\text{ is a connected component of }(\tilde{J}+\alpha)\setminus \tilde{J'},\alpha_z<0\}.\]

Next, we define the set $B^{J,J'}_n$ as an extension of $A^{J,J'}_n$ by
\begin{align*}    
B^{J,J'}_n=:\{(U,\alpha)\in\mathcal{P}(\mathbb{Z}\times\mathbb{N})\times\mathbb{Z}^2:U\text{ is a connected component of } ( \tilde{J}+\alpha)\setminus\tilde{J'},\alpha_z<0\}.\end{align*}
Note that $B^{J,J'}_n$ corresponds to the not necessarily bounded components of $(\tilde{J}+\alpha)\setminus\tilde{J'}$ with $\alpha_z<0$, which are contained entirely in the upper-half-plane.


    
\end{definition}

\begin{proof}[Proof of Lemma \ref{Lem:inequalityL1L2}]
 
First, we biject $A^{J,J'}_p$ with $\mathbb{N}^2\setminus\tilde{J'}$, which in turn corresponds to the natural basis of $S/J'$; hence $A^{J,J'}_p$ will be of
size $l'$. To construct such a bijection,  for $(U,\alpha)\in A^{J,J'}_p$, we choose  $\gamma={(\gamma_y,\gamma_z})$ in $\tilde{J}$ with maximal $z$-coordinate and $\gamma+\alpha\in U$. If there is more than one option, we take the one with minimal $y$-coordinate. Similarly, we have that $\gamma_y+\alpha_y$ is the smallest $y$-coordinate of any element of $U$.

Define
\begin{align*}
    f:A^{J,J'}_p&\rightarrow \mathbb{N}^2\setminus\tilde{J'}\\
    (U,\alpha)&\mapsto(\gamma_y+\alpha_y,\alpha_z)
\end{align*}




\emph{Claim 1.} $f$ is a well-defined bijection.
\begin{proof}[Proof of Claim 1] We prove this in three steps:


\begin{enumerate}
    \item $f$ is well-defined: $\gamma_y+\alpha_y\in\BN$, and $\alpha_z\in \BN$ are clear. If $f(U,\alpha)\in\tilde{J'}$, then since $\gamma_z\geq0$, we must have $\alpha+\gamma\in\tilde{J'}$ and $\alpha+\gamma\in U$; contradiction.
    


     \item $f$ is a surjection: Given an element $\beta=(\beta_y,\beta_z)$ in $\mathbb{N}^2\setminus\tilde{J'}$, we have $\beta_z\geq 0$. Hence, we can define $\alpha_z:=\beta_z$. Now, we define $\alpha_y$ in the following way. Take $\alpha_y$ to be maximal with the property that there are elements with $y$-coordinate equal to $\beta_y$ in $(\tilde{J}+\alpha)\setminus\tilde{J'}$. Now, take $U$ to be the unique component of $(\tilde{J}+\alpha)\setminus\tilde{J'}$ which contains the elements with  $y$-coordinate equal to $\beta_y$. If $U$ contains elements with $y$-coordinate less than $\beta_y$, we arrive at a contradiction because we would have two points of coordinates $(\beta_y,s),(\beta_y-1,s)\in U$ so we can increase $\alpha_y$ by at least 1. Note that since $\alpha_z\geq 0$, $U$ is restricted to the upper-half-plane. On the other hand, since  $\beta_y\geq 0$, $U$ is also restricted to the right-half-plane. Hence, $U$ is bounded, and we have $f(U,\alpha)=\beta$.

     \item $f$ is an injection: Suppose $\beta=(\gamma_y+\alpha_y,\alpha_z)=f(U,\alpha)=f(\bar{U},\bar{\alpha})=(\bar{\gamma}_y+\bar{\alpha}_y,\bar{\alpha}_z)$, for some $\bar{\alpha}=(\bar{\alpha}_y,\bar{\alpha}_z)$, $\bar{\gamma}=(\bar{\gamma}_y,\bar{\gamma}_z)$ and $\bar{U}$, defined similarly as $\alpha$, $\gamma$ and $U$. It is clear that $\alpha_z=\bar{\alpha}_z$. For the sake of contradiction, suppose $\alpha_y<\bar{\alpha}_y$; then $\gamma_y>\bar{\gamma}_y$. Now $\bar{\gamma}+\bar{\alpha}=\bar{\gamma}+\alpha+r\mathbf{e}_y\notin\tilde{J}'$ where $r=\bar{\alpha}_y-\alpha_y>0$. We deduce that $\bar{\gamma}+\alpha\in U$, which contradicts the definition of $\gamma$. Hence, $\alpha=\bar{\alpha}$. Finally, $U$ and $\bar{U}$ are the unique connected component of $(\tilde{J}+\alpha)\setminus\tilde{J}'$, which contain elements with $y$-coordinate equal to $\beta_y$. Thus $U=\bar{U}$, and we have an injection as required.
\end{enumerate}
\end{proof}
 

We can similarly biject $B^{J,J'}_n$ with $\mathbb{N}^2\setminus\tilde{J}$. For  this purpose, associated to $(U,\alpha)\in B^{J,J'}_n$, we choose $\gamma=(\gamma_y,\gamma_z)$ in $\tilde{J}$ such that $\gamma+\alpha\in U$, and $\gamma$ has maximal $y$-coordinate. If there are multiple options, we take the one with minimal $z$-coordinate. 

Define
\begin{align*}
    g:B^{J,J'}_n&\rightarrow \mathbb{N}^2\setminus\tilde{J}\\
    (U,\alpha)&\mapsto (\gamma_y,-\alpha_z-1)
\end{align*}

\emph{Claim 2.} $g$ is a well-defined bijection.

\begin{proof}[Proof of Claim 2] Again, we proceed with the proof in three steps:

\begin{enumerate}
    \item $g$ is well-defined: It is clear that $\gamma_y \in \BN$ and $-\alpha_z-1\in\BN$. Since $U$ is in the upper-half-plane, we have $\gamma_z+\alpha_z+1> 0$. Therefore, $\gamma_z>-\alpha_z-1$, and so $f(U,\alpha)\notin\tilde{J}$.
    \item $g$ is a surjection: Given an element $\beta$ in $\mathbb{N}^2\setminus\tilde{J}$, we can find $\gamma=(\gamma_y,\gamma_z)\in J$ such that $\gamma_y=\beta_y$ and $\gamma_z$ is minimal. We set $\alpha_z:=-\beta_z-1$. Then, we can find the unique $\lambda$ such that $(\lambda,\gamma_z+\alpha_z)\notin\tilde{J'}$ but $(\lambda+1,\gamma_z+\alpha_z)\in\tilde{J'}$. Then, we can uniquely define $\alpha_y:=\lambda-\gamma_y$ and $\alpha=(\alpha_y,\alpha_z)$. Choosing $U$ to be the connected component of $(\tilde{J}+\alpha)\setminus\tilde{J}'$ containing $\alpha+\gamma$, then as before, $\gamma_z+\alpha_z$ is the minimal $z$-coordinate of an element of $U$. So, this gives $f(U,\alpha)=\beta$.
    \item $g$ is an injection: This part is similar to the injection of $f$ as above.
\end{enumerate}
    
\end{proof}


Combining Claim 1 and Claim 2, we have
\begin{align*}
    \mathrm{hom}(J,S/J')= l+l'-\mathrm{Card}( B^{J,J'}_n\setminus A^{J,J'}_n),
    \end{align*}
which implies the desired inequality.
\end{proof}

\begin{remark}
    The bound \eqref{eq: Ineq} is related to  \cite[Proposition 4.1]{Ramkumar-Sammartano}. We prove \eqref{eq: Ineq}  via a different method, which is purely combinatorical. This approach allows us to improve the bounds by subtracting further terms, which we can evaluate easily in order to eventually prove the original conjecture.
\end{remark}






\begin{corollary} \label{Cor:An'An}
    For $J$ and $J'$ as above, we have
    \begin{align}
        \mathrm{hom}(J,S/J')=l+l'-\mathrm{Card}( B^{J,J'}_n\setminus A^{J,J'}_n)
    \end{align}
\end{corollary}
\begin{proof}
    This is immediate from the proof of Lemma \ref{Lem:inequalityL1L2}.
\end{proof}

Moreover, Lemma \ref{Lem:inequalityL1L2} immediately implies the following well-known result about the Hilbert scheme of $\BA^2$.
\begin{corollary}
     $\Hilb^l(\mathbb{A}^2)$ is smooth.
\end{corollary}
\begin{proof}
    Taking $I=J=J'$ we get the inequality $T(I)\leq 2l$ which suffices (since $2l \leq T(I)$) to deduce the smoothness of $\Hilb^l(\mathbb{A}^2)$.
\end{proof}
\begin{definition}
    Let $J$ be a $0$-dimensional monomial ideal in $\BC[y,z]$ we define its \emph{height} to be $h$ such that $z^h$ is a minimal generator.
\end{definition}
\begin{definition}
    Let $I_0,I_1$ be two $0$-dimensional monomial ideals in $\BC[y,z]$. We say that $I_0$ is \emph{taller than} $I_1$ (or equivalently, $I_1$ is \emph{shorter than} $I_0$) if the height of $I_0$ is greater than that of $I_1$.
\end{definition}

\begin{definition} Let $I_0,I_1$ be two $0$-dimensional  monomial ideals in $\BC[y,z]$ such that $I_1$ has height $h$. We define $t(I_0,I_1)$ to be the size of the set $\{\beta\in\mathbb{N}^2\setminus\tilde{I_0}:\beta_z\geq h\}$.

    Note that if $I_1$ is taller than $I_0$, then $t(I_0,I_1)=0$.
\end{definition}
 




\begin{lemma}\label{Cor:tLessGhost}
 For $J$ and $J'$ as above, we have $t(J,J') \leq \mathrm{Card}(B^{J,J'}_n\setminus A^{J,J'}_n)$.
\end{lemma}



\begin{proof}

   Let $h$ be the height of $J'$. Each point $\beta=(\beta_y,\beta_z)$ counted by $t(J,J')$ can be mapped to $(-1,h-1)$ and these correspond to distinct elements of $B^{J,J'}_n\setminus A^{J,J'}_n$, although not necessarily all of them. Since the value of $\alpha=(\alpha_y,\alpha_z)$ used is $(-1,h-1)-\beta=(-1-\beta_y,h-1-\beta_z)$, we have $\alpha_z=h-1-\beta_z<0$, since by definition, $\beta_z\geq h$. Hence, it suffices to check that there is a connected component $U$ of $(\tilde{J}+\alpha)\setminus\tilde{J'}$, which is unbounded but only inside the upper-half-plane. First, we note that $\alpha_y=-1-\beta_y<0$ so there is certainly an unbounded component $U$ in the upper-left quadrant. However, by the definition of $\beta$ and $h$, we have $\beta\notin\tilde{J}$ and $(0,h)\in\tilde{J'}$, respectively. The former implies that there is no $(\Fy,\Fz)$ in $(\tilde{J}+\alpha)$ such that $\Fy\leq-1,\Fz\leq h-1$, as $\beta+\alpha=(-1,h-1)$. The latter implies that for any $(\Fy,\Fz)$ such that $\Fy\geq 0,\Fz\geq h$, we have $(\Fy,\Fz)\in\tilde{J'}$. Combining these implies that $U\subset (-\infty,-1]\times[h,\infty)$; thus, this corresponds to an element of $B^{J,J'}_n\setminus A^{J,J'}_n$. Note that for different values of $\beta$ we get different values of $\alpha$, so we get $t(J,J')$ distinct elements of $B^{J,J'}_n$.
\end{proof}

The following statement is immediate from Corollary \ref{Cor:An'An} and Lemma \ref{Cor:tLessGhost}. 

\begin{corollary}\label{cor:ineq}
    For $J$ and $J'$ as above, we have $\mathrm{hom}(J,S/J')\leq l+l'-t(J,J')$.
\end{corollary}

\section*{Retained context: the preparatory lemma and corollary (paper \S2.1)}
First, in order to prove Theorem \ref{thm:upperbound}, we need a preparatory lemma and its immediate corollary concerning a lower bound for $t(I_i,I_j)$.
\begin{lemma}\label{Lem:bound}
    Let $I$ be a $0$-dimensional Borel-fixed ideal in $\BC[x,y,z]$, and let $\alpha=(\alpha_x,\alpha_y,\alpha_z)$ be a vector in $\BZ^3$ such that $\alpha_x<0$ and $\alpha_y,\alpha_z\geq0$ with $\alpha_x+\alpha_y+\alpha_z\leq-1$. Then, for any minimal generator $\gamma$ in $\tilde{I}$, we have $\gamma+\alpha\notin\tilde{I}$.
\end{lemma}
\begin{proof}
    If $\gamma+\alpha\in\tilde{I}$, then by Borel-fixedness and that $\alpha_x<0$, we deduce that $\gamma-\mathbf{e}_x\in\tilde{I}$, which contradicts the fact that $\gamma$ corresponds to a generator of $I$. 
\end{proof}

\begin{corollary}\label{cor:tlowerbound}
    For $I$ a 0-dimensional Borel-fixed ideal in $\BC[x,y,z]$, with $I=\bigoplus_ix^iI_i$ then
    \[\binom{j-i+1}{2} \leq t(I_i,I_j).\]
\end{corollary}

We now prove Theorem \ref{thm:upperbound}, following a similar argument as in \cite{Ramkumar-Sammartano}, but using our slightly improved inequalities to get a stronger bound which we can use to deduce, for example, the Briançon-Iarrobino Conjecture for $N=3$.

\section*{The authors' complete proof of Theorem A (paper \S2.1)}
\begin{proof}[Proof of Theorem \ref{thm:upperbound}]
Recall that we have the decomposition $I=\bigoplus_ix^iI_i$. Any element of the tangent space $\mathrm{Hom}(I,R/I)$ can restrict to a non-zero element of $\mathrm{Hom}(I_i,S/I_j)$ for some $i,j$, and for basis elements, the restriction determines the original map. 

\emph{Claim.} Let $i>j$. Then, maps in $\mathrm{Hom}(I_i,S/I_j)$ with $\alpha$ such that $\alpha_y,\alpha_z\geq0$ and $\alpha_y+\alpha_z\leq i-j-2$ correspond to a single bounded connected component, but cannot extend to elements of the tangent space $\Hom(I,R/I)$, i.e., they are zero vectors.

\begin{proof}[Proof of Claim]
    Let $\alpha_x=j-i$, then by Lemma \ref{Lem:bound} for any minimal generator $\gamma\in\tilde{I}$ of $I$, we have
   \begin{align}
\gamma+\alpha+\mathbf{e}_y \notin \tilde{I},\label{1}\\
\gamma+\alpha+\mathbf{e}_x\notin\tilde{I}\label{2}.
   \end{align}
   
   Therefore, \eqref{1} implies that for generators $\gamma$ and $\gamma'$ whose $y$-coordinates differ by $1$, we have that $\gamma+\alpha$ and $\gamma'+\alpha$ are in the same component; hence, we have a single connected component in $(I_i+(\alpha_y,\alpha_z))\setminus I_j$. Also, we note that since $\alpha_y,\alpha_z\geq0$, the component is bounded.

    Similarly, \eqref{2} implies that there is only one connected component in $\tilde{I}_\alpha$, since for generators $\gamma$ and $\gamma'$ whose $z$-coordinates differ by 1, we have that  $\gamma+\alpha$ and $\gamma'+\alpha$ are in the same component. However, since $\alpha_x<0$, the component is unbounded.
\end{proof}
For fixed $i,j$, the number of maps in the claim above is $\binom{i-j}{2}$.

Also, we have 

\begin{align}\label{eq:doublebinomial}
    \sum_{j=0}^{m_1-1}\Bigg(\sum_{i=j+1}^{m_1}\binom{i-j}{2}\Bigg)= \sum_{j=0}^{m_1-1}{m_1-j+1\choose 3}={m_1+2 \choose 4}.
\end{align}

Thus, by the claim above, Corollary \ref{cor:ineq}, Corollary  \ref{cor:tlowerbound} and using \eqref{eq:doublebinomial} twice, it is straightforward to show
\begin{align*}
    T(I)&=\mathrm{hom}(I,R/I)\\
    &\leq\sum_{i=0}^{m_1}\sum_{j=0}^{m_1-1}\mathrm{hom}(I_i,S/I_j)-\sum_{j=0}^{m_1-1}\sum_{i=j+1}^{m_1}\binom{i-j}{2}\\
    &\leq\sum_{i=0}^{m_1}\sum_{j=0}^{m_1-1}\big(l_i+l_j-t(I_i,I_j)\big)-\binom{m_1+2}{4}\\
    &=\sum_{i=0}^{m_1}m_1l_i+\sum_{j=0}^{m_1-1}(m_1+1)l_j-\sum_{i=0}^{m_1}\sum_{j=0}^{m_1-1}t(I_i,I_j)-\binom{m_1+2}{4}\\
    &\leq(2m_1+1)l-\sum_{i=0}^{m_1}\sum_{j=i+1}^{m_1-1}\binom{j-i+1}{2}-\binom{m_1+2}{4}\\
    &=(2m_1+1)l-2\binom{m_1+2}{4},
\end{align*}
as required.
\end{proof}

\begin{thebibliography}{99}
\bibitem[37]{Ramkumar-Sammartano} R.~Ramkumar and A.~Sammartano, {\em {On the tangent space to the Hilbert scheme of points in $\mathbb{P}^3$}}, Transactions of the American Mathematical Society {\bf 375} (2022), 6179--6203.
\end{thebibliography}
\end{document}
